\documentclass[11pt]{article}
% Shared local declarations. Document roots use the repository's business-admin family.
\usepackage{business-admin}
\usepackage{amsmath,amssymb,tabularx,longtable,array,booktabs,enumitem}
\usepackage{ragged2e}
\newcolumntype{Y}{>{\RaggedRight\arraybackslash}X}
\newcommand{\SuiteMeta}[2]{\def\DocTitle{#1}\def\DocSubtitle{#2}\def\DocVersion{1.0}\def\DocStatus{Educational and operational reference}\def\DocOwner{Parliamentary Procedure and Governance Suite}\def\DocSummary{Independently authored reference and working materials, based on RONR 12th edition.}}
\newcommand{\SuiteOpen}{\makedoctitle\vspace{-0.5em}\noindent\rule{\linewidth}{0.4pt}\par\medskip}
\newcommand{\RonrAnchor}[1]{\par\smallskip{\footnotesize\textbf{Primary reading:} Robert's Rules of Order Newly Revised, 12th ed., #1. Original synthesis; not official interpretation.}\par}
\newcommand{\WriteLine}{\par\noindent\rule{\linewidth}{0.35pt}\par\vspace{0.5cm}}
\newcommand{\WriteLines}[1]{\par\noindent\rule{\linewidth}{0.35pt}\par\vspace{0.35cm}\noindent\rule{\linewidth}{0.35pt}\par\vspace{0.35cm}\noindent\rule{\linewidth}{0.35pt}\par\vspace{0.25cm}}
\newcommand{\Field}[1]{\par\medskip\noindent\textbf{#1}\par\WriteLines{3}}
\newcommand{\GovCheckbox}{\(\square\)}
\setlength{\emergencystretch}{3em}

\SuiteMeta{Section 12: Amend}{Study guide | Part VI — Subsidiary Motions}
\begin{document}
\SuiteOpen
\RonrAnchor{Section 12}
\section{Study outcomes}
\begin{itemize}[leftmargin=*,itemsep=4pt]
\item Explain the procedural purpose of amend and locate the governing rules in Section 12.
\item Track first- and second-degree amendments; maintain grammatical and substantive coherence; use filling blanks when appropriate.
\item Produce a defensible meeting record showing the governing authority, procedural decision, and result.
\end{itemize}

\section{Conceptual synthesis}
Track first- and second-degree amendments; maintain grammatical and substantive coherence; use filling blanks when appropriate.

This topic operates within the assembly’s adopted hierarchy of authority: applicable law and charter, bylaws, special rules of order, the parliamentary authority where applicable, standing rules, and valid custom. The exact procedure depends on the kind of body, the status of pending business, notice, and the governing documents. The 12th edition should be consulted for conditions and exceptions rather than treating a summary matrix as binding law.
\section{Decision procedure}
\begin{enumerate}[leftmargin=*,itemsep=4pt]
\item Classify the body and confirm who holds voting and participation rights.
\item Find the controlling bylaw, special rule or cited RONR paragraph; record any conflict.
\item Identify the pending question, the member’s procedural request and the relevant time or notice restriction.
\item Check the requirement for a second, debate, amendment, threshold and announcement where relevant.
\item Preserve the operative wording, decision, and any required supporting evidence.
\end{enumerate}

\section{Laboratory: scenario and evidence}
Process insertion, striking out, and substitution without losing the current text.

\textbf{Required artifacts:} procedural timeline; controlling authority citation; chair ruling or vote announcement; concise minutes extract; one alternative branch with a different governing rule.

\section{Failure mode and correction}
\textbf{Risk:} Voting on the main motion before pending amendments are resolved.

\textbf{Control:} pause the workflow, identify the authority and pending question, explain the condition to the assembly, and document the resolution without retroactively rewriting an announced result.

\section{Knowledge check}
\begin{enumerate}[leftmargin=*,itemsep=4pt]
\item How many degrees of amendment may be pending?
\item How would the resolution change if a higher-priority bylaw differed from the usual rule?
\item What minimum evidence would allow another secretary to reconstruct the result?
\end{enumerate}

\section{Assessment standard}
Credit requires an accurate citation to RONR (12th ed.) Section 12, a correct classification of the body or motion, a reproducible action sequence, and a meeting record with no invented facts.

\section{Adjacent reading}
Use the table of contents and the Charts, Tables, and Lists at the end of RONR for related rules. Cross-check Sections 2, 4, 7, 40, 44, 45, and 48 when a case involves validity, voting, or records.
\end{document}
